\begin{tikzpicture}[
  x=1cm,y=1cm,
  participant/.style={draw=blue!45!black,fill=blue!4,rounded corners=3pt,
    minimum width=7.8cm,minimum height=1.05cm,align=center,font=\small},
  message/.style={-{Latex[length=2.7mm]},line width=.9pt,draw=black!75},
  reply/.style={message,dashed},
  msglabel/.style={font=\small,fill=white,inner sep=3pt,align=center},
  note/.style={anchor=north west,text width=14.3cm,align=left,font=\small,
    execute at begin node=\raggedright},
  guard/.style={anchor=north west,font=\small\bfseries,fill=white,inner sep=3pt},
  lifeline/.style={draw=black!35,dashed,line width=.7pt}
]
\node[anchor=north west,font=\LARGE\bfseries] at (0,0)
  {Login: frontend--backend communication};
\node[anchor=north west,text width=38cm,font=\small] at (0,-.85)
  {Time runs downward. Solid arrows are calls; dashed arrows are responses.
   The framed outcomes summarize alternative replies, not consecutive events;
   early failures stop immediately at the indicated step. Blue source labels are clickable.};

\node[participant] (frontend) at (4,-2.25) {
  \textbf{Frontend}\\
  \source{host/assets/app.js}{JavaScript UI / fetch client}
};
\node[participant] (backend) at (19,-2.25) {
  \textbf{Backend}\\
  \source{host/main.rs}{Rust HTTP boundary} +
  \source{theories/Api.v}{Rocq service} + provider interfaces
};
\draw[lifeline] (4,-2.8) -- (4,-24.1);
\draw[lifeline] (19,-2.8) -- (19,-24.1);

\draw[message] (4,-4) -- node[msglabel,above] {
  \code{POST /api/login}\quad\code{Content-Type: application/json}\\
  \code{{"email": "...", "password": "..."}}
} (19,-4);
\node[anchor=north west,text width=9.6cm,font=\small,align=left] at (6,-4.8) {
  The user submits email and password.\par
  Passwords are not put into URLs or browser storage, and are not echoed in responses.
};

\draw[message] (19,-5) -- (20.5,-5) -- (20.5,-5.45) -- (19,-5.45);
\node[note] at (22,-4.85) {
  \textbf{1. Check the HTTP envelope.}\par
  \source{host/main.rs}{Host/Origin, method, content type and 16 KiB limit}.
  Boundary rejection ends the request without a credential lookup.
};

\draw[message] (19,-6.65) -- (20.5,-6.65) -- (20.5,-7.1) -- (19,-7.1);
\node[note] at (22,-6.5) {
  \textbf{2. Enter the component and read the clock.}\par
  \source{component/wit/service.wit}{\code{login(config, body:list<u8>) -> response}}.
  \source{component/adapters/service.wat}{The wrapper reads WASI wall time once}; pure Rocq receives Unix seconds.
};

\draw[message] (19,-8.3) -- (20.5,-8.3) -- (20.5,-8.75) -- (19,-8.75);
\node[note] at (22,-8.15) {
  \textbf{3. Parse and validate the request.}\par
  \source{theories/Api.v}{\code{parse_request RtLogin body}} requires email/password strings.
  Malformed JSON or missing/wrongly typed fields returns \code{400}.
};

\draw[message] (19,-9.95) -- (20.5,-9.95) -- (20.5,-10.4) -- (19,-10.4);
\node[note] at (22,-9.8) {
  \textbf{4. Fetch and decode the account.}\par
  \source{component/adapters/store.wat}{WASI keyvalue bucket get(email)};
  \source{theories/Store.v}{Rocq decodes the record and checks its email}.
  Store refusal returns \code{500}. Missing/corrupt/mismatched records select dummy credentials.
};

\draw[message] (19,-11.65) -- (20.5,-11.65) -- (20.5,-12.1) -- (19,-12.1);
\node[note] at (22,-11.5) {
  \textbf{5. Verify the password inside the crypto provider.}\par
  \source{component/adapters/argon2.wat}{Argon2 verification} receives the real or dummy salt, parameters and tag.
  \source{theories/Api.v\#L104}{\code{login_check}} hashes even when the user is absent.
  This is not a constant-time theorem.
};

% Mutually exclusive replies. Early exits refer back to the numbered steps.
\draw[draw=black!45,rounded corners=2pt] (0,-13.5) rectangle (38,-24);
\node[guard] at (.25,-13.38) {alt: exactly one outcome};
\node[guard] at (.4,-14.05) {Found user + valid password + successful signing};

\draw[message] (19,-14.8) -- (20.5,-14.8) -- (20.5,-15.25) -- (19,-15.25);
\node[note] at (22,-14.65) {
  \textbf{6. Mint the token.}\par
  \source{theories/Jwt.v}{Pinned EdDSA header; claims \code{sub,iss,iat,exp}}.
  \source{component/adapters/ed25519.wat}{Sign with the opaque key handle}.
  Password and private key are not included in the JWT.
};
\draw[reply] (19,-16.65) -- node[msglabel,above] {
  \code{200 OK}\quad\code{{"token": "header.payload.signature"}}
} (4,-16.65);

\draw[draw=black!30,dashed] (0,-17.45) -- (38,-17.45);
\node[guard] at (.4,-17.6) {Missing account, wrong password, or Argon2 refusal};
\draw[reply] (19,-18.85) -- node[msglabel,above] {
  \code{401 Unauthorized}\quad\code{{"error": "unauthorized"}}
} (4,-18.85);
\node[note] at (22,-18.1) {
  No signing call. An absent user is rejected even if dummy verification says true.
  The response does not distinguish missing accounts from wrong passwords.
};

\draw[draw=black!30,dashed] (0,-19.5) -- (38,-19.5);
\node[guard] at (.4,-19.65) {Password accepted, but signing provider refuses};
\draw[reply] (19,-20.8) -- node[msglabel,above] {
  \code{500 Internal Server Error}\quad\code{{"error": "internal error"}}
} (4,-20.8);
\node[note] at (22,-20.05) {
  A failed token-signing operation is an internal failure, not a credential rejection.
};

\draw[draw=black!30,dashed] (0,-21.5) -- (38,-21.5);
\node[guard] at (.4,-21.65) {Early exit from request parsing or storage};
\draw[reply] (19,-22.95) -- node[msglabel,above] {
  \code{400 Bad Request}\quad or\quad\code{500 Internal Server Error}
} (4,-22.95);
\node[note] at (22,-22) {
  \code{400}: malformed request (step 3).\par
  \code{500}: store unavailable (step 4), before hashing.\par
  HTTP-envelope errors (step 1) may instead be \code{403/405/413/415}.
};

\node[anchor=north west,text width=38cm,font=\small] at (0,-24.65) {
  \textbf{After the reply:} the frontend displays success or error; on success it
  copies the returned token into the validation field. Token validation is a
  separate request, not part of login. Login does not write account records or
  send email. \source{host/assets/app.js}{UI implementation}\quad
  \source{tests/http.rs}{HTTP flow tests}\quad
  \source{theories/Proofs.v\#L99}{No-token-without-check statement}.
  This sequence describes behavior; it is not an end-to-end security proof.
};
\end{tikzpicture}
